\begin{table}[!htb]
\centering
\small
\begin{tabular}{lrrrr}
\toprule
\textbf{Shock} & \textbf{Min} & \textbf{Max} & \textbf{Mean} & \textbf{Median} \\
\midrule
$\text{shock\_wage}$    & $-4.2\%$ (NE) & $+0.6\%$ (ND) & $-1.0\%$ & $-0.9\%$ \\
$\text{shock\_amenity}$ & $-3.5\%$ (ME) & $+7.8\%$ (FL) & $+1.0\%$ & $+0.6\%$ \\
$\text{shock\_housing}$ & $-0.0\%$ (WA) & $+0.4\%$ (AR) & $+0.2\%$ & $+0.2\%$ \\
\bottomrule
\end{tabular}
\caption{Summary statistics of the three composite climate shocks across 51 states (DC included). Signs: negative wage $=$ loss; positive amenity $=$ welfare deterioration (higher mortality); positive housing $=$ higher housing cost. Note that $\text{shock\_amenity}$ is reported in consumption-equivalent percentage terms (the absolute util shock divided by the CARA parameter $\gamma = 3$).}
\label{tab:shock_magnitudes}
\end{table}